\SourceLeaf{538}\par\medskip\Needspace{5\baselineskip}\textit{Nachlass.}
\begingroup\fontsize{10}{13}\selectfont
\begin{longtable}{@{}p{.86\linewidth}@{\quad}r@{}}
\multicolumn{2}{@{}l}{Analysis residuorum:}\\[3pt]
Caput sextum. Pars prior. Solutio congruentiae $x^m-1\equiv0$ .\,.\,.\,.\,.\,.\,.\,.\,. & Seite 199 \\[3pt]
Caput octavum. Disquisitiones generales de congruentiis .\,.\,.\,.\,.\,.\,.\,.\,.\,. & — 212 \\[3pt]
Disquisitionum circa aequationes puras ulterior evolutio .\,.\,.\,.\,.\,.\,.\,.\,.\,.\,.\,. & — 243 \\[3pt]
Démonstration de quelques théorèmes concernants les périodes des classes des formes binaires du second degré .\,.\,.\,.\,.\,.\,.\,.\,.\,.\,.\,.\,.\,.\,.\,.\,.\,.\,.\,.\,.\,.\,.\,. & — 266 \\[3pt]
De nexu inter multitudinem classium in quas formae binariae secundi gradus distribuuntur earumque determinantem. I. II\,.\,. X .\,.\,.\,.\,.\,.\,.\,.\,.\,.\,.\,.\,.\,.\,.\,.\,. & — 269 \\[3pt]
Geometrische Seite der ternären Formen .\,.\,.\,.\,.\,.\,.\,.\,.\,.\,.\,.\,.\,.\,.\,.\,. & — 305 \\[3pt]
Zur Theorie der biquadratischen Reste. I.\,.\,. VI .\,.\,.\,.\,.\,.\,.\,.\,.\,.\,.\,.\,.\,. & — 313 \\[3pt]
Zur Theorie der complexen Zahlen. I.\,.\,. VI .\,.\,.\,.\,.\,.\,.\,.\,.\,.\,.\,.\,.\,.\,. & — 387 \\[3pt]
Tafel des quadratischen Characters der Primzahlen .\,.\,.\,.\,.\,.\,.\,.\,.\,.\,.\,.\,.\,. & — 399 \\[3pt]
Tafel zur Verwandlung gemeiner Brüche in Decimalbrüche .\,.\,.\,.\,.\,.\,.\,.\,.\,.\,. & — 411 \\[3pt]
Tafel der Frequenz der Primzahlen .\,.\,.\,.\,.\,.\,.\,.\,.\,.\,.\,.\,.\,.\,.\,.\,.\,.\,.\,. & — 435 \\[3pt]
Tafel der Anzahl der Classen binärer quadratischer Formen .\,.\,.\,.\,.\,.\,.\,.\,.\,.\,. & — 449 \\[3pt]
Tafel zur Cyklotechnie .\,.\,.\,.\,.\,.\,.\,.\,.\,.\,.\,.\,.\,.\,.\,.\,.\,.\,.\,.\,.\,.\,.\,. & — 477 \\[3pt]
Circuli quadratura nova .\,.\,.\,.\,.\,.\,.\,.\,.\,.\,.\,.\,.\,.\,.\,.\,-\,.\,.\,.\,.\,.\,.\,. & — . 497 \\[3pt]
Zur Berechnung der Logarithmen .\,.\,.\,.\,.\,.\,.\,.\,.\,.\,.\,.\,.\,.\,.\,.\,.\,.\,.\,. & — 501 \\[3pt]
Quadratorum myrias prima .\,.\,.\,.\,.\,.\,.\,.\,.\,.\,.\,.\,.\,.\,.\,.\,.\,.\,.\,.\,.\,.\,. & — 504 \\[3pt]
Indices der Primzahlen im höhern Zahlenreiche .\,.\,.\,.\,.\,.\,.\,.\,.\,,\,.\,.\,.\,.\,. & — 506 \\[3pt]
Hülfstafel zur Auflösung der Gleichung $A=fxx+gyy$ .\,.\,.\,.\,.\,.\,.\,.\,.\,.\,.\,. & — 509 \\[3pt]
Sectio octava. Quarundam disquisitionum ad circuli sectionem pertinentium uberior consideratio.  & — 510 \\[3pt]
Briefe von \textsc{Gauss} an \textsc{Dirichlet} .\,.\,.\,.\,.\,.\,.\,.\,.\,.\,.\,.\,.\,.\,.\,.\kern.6em.\,.\,.\,. & — 514 \\[3pt]
\end{longtable}\endgroup
\begingroup\fontsize{10}{13}\selectfont
\begin{longtable}{@{}p{.86\linewidth}@{\quad}r@{}}
\textit{Bemerkungen} .\,.\,.\,.\,.\,.\,.\,.\,.\,.\,.\,.\,.\,.\,.\,.\,.\,.\,.\,.\,.\,.\,.\,.\,.\,.\,.\,. & — 519 \\[3pt]
\end{longtable}\endgroup
\BookRule{.25\textwidth}
\begin{center}
GÖTTINGEN,\par\smallskip
DRUCK DER DIETRICHSCHEN UNIVERSITÄTS-BUCHDRUCKEREI.\par\smallskip
W. FR. KAESTNER.
\end{center}
